\documentclass[conference,a4paper,10pt]{IEEEtran}

%% --- движок и шрифты ---
\usepackage{fontspec}
\usepackage{polyglossia}
\setdefaultlanguage{russian}
\setotherlanguage{english}
\setmainfont{Linux Libertine O}
\setsansfont{Linux Biolinum O}
\setmonofont{DejaVu Sans Mono}[Scale=0.82]

%% --- стандартные пакеты ---
\usepackage{amsmath, amssymb}
\usepackage{graphicx}
\usepackage{hyperref}
\usepackage{booktabs}
\usepackage{listings}
\usepackage{xcolor}
\usepackage{caption}
\usepackage{cite}

\lstset{
  basicstyle=\ttfamily\footnotesize,
  breaklines=true,
  frame=single,
  columns=fullflexible,
  keepspaces=true,
  showstringspaces=false,
  literate={а}{{\selectfont а}}1 {б}{{\selectfont б}}1 {в}{{\selectfont в}}1
           {г}{{\selectfont г}}1 {д}{{\selectfont д}}1 {е}{{\selectfont е}}1
           {ж}{{\selectfont ж}}1 {з}{{\selectfont з}}1 {и}{{\selectfont и}}1
           {й}{{\selectfont й}}1 {к}{{\selectfont к}}1 {л}{{\selectfont л}}1
           {м}{{\selectfont м}}1 {н}{{\selectfont н}}1 {о}{{\selectfont о}}1
           {п}{{\selectfont п}}1 {р}{{\selectfont р}}1 {с}{{\selectfont с}}1
           {т}{{\selectfont т}}1 {у}{{\selectfont у}}1 {ф}{{\selectfont ф}}1
           {х}{{\selectfont х}}1 {ц}{{\selectfont ц}}1 {ч}{{\selectfont ч}}1
           {ш}{{\selectfont ш}}1 {щ}{{\selectfont щ}}1 {ъ}{{\selectfont ъ}}1
           {ы}{{\selectfont ы}}1 {ь}{{\selectfont ь}}1 {э}{{\selectfont э}}1
           {ю}{{\selectfont ю}}1 {я}{{\selectfont я}}1
           {А}{{\selectfont А}}1 {Б}{{\selectfont Б}}1 {В}{{\selectfont В}}1
           {Г}{{\selectfont Г}}1 {Д}{{\selectfont Д}}1 {Е}{{\selectfont Е}}1
           {Ж}{{\selectfont Ж}}1 {З}{{\selectfont З}}1 {И}{{\selectfont И}}1
           {Й}{{\selectfont Й}}1 {К}{{\selectfont К}}1 {Л}{{\selectfont Л}}1
           {М}{{\selectfont М}}1 {Н}{{\selectfont Н}}1 {О}{{\selectfont О}}1
           {П}{{\selectfont П}}1 {Р}{{\selectfont Р}}1 {С}{{\selectfont С}}1
           {Т}{{\selectfont Т}}1 {У}{{\selectfont У}}1 {Ф}{{\selectfont Ф}}1
           {Х}{{\selectfont Х}}1 {Ц}{{\selectfont Ц}}1 {Ч}{{\selectfont Ч}}1
           {Ш}{{\selectfont Ш}}1 {Щ}{{\selectfont Щ}}1
           {ё}{{\selectfont ё}}1 {Ё}{{\selectfont Ё}}1
}

\begin{document}

\title{Минималистичная реализация эффективного хранения JSON в произвольных реляционных СУБД через декомпозицию путей в side-таблицы}

\author{\IEEEauthorblockN{Малько Е.\,А.}
\IEEEauthorblockA{Университет ИТМО\\
Санкт-Петербург, Россия\\
egormalko1@outlook.com}}

\maketitle

\begin{abstract}
Настоящая работа адресована \emph{разработчикам СУБД}, перед которыми
стоит задача добавить в свой движок эффективную поддержку
полу\-струк\-турированных данных (JSON). За основу взят подход
ClickHouse, в котором каждый встретившийся в документах JSON-путь
раскладывается в отдельную колонку. Этот подход обеспечивает
почти-нативную колоночную производительность, однако требует
эффективного хранения колонок с очень большим числом NULL-значений;
в ClickHouse это решено на физическом уровне специализированными
структурами хранения. В произвольной реляционной СУБД такой
инфраструктуры может не быть, и добавлять её ради поддержки JSON
дорого. Мы показываем, что ту же задачу можно решить уже
имеющимися в практически любом реляционном движке механизмами:
обычными таблицами, хеш-индексами, MVCC и стандартными физическими
операторами scan/join. В предложенной системе данные раскладываются
не по колонкам, а по отдельным таблицам, в которых лежат только
non-null значения, а оригинальная структура восстанавливается с
помощью оператора join. Кроме того, рассматривается вопрос корректного
поведения при параллельных запросах, а также некоторые эвристики,
призванные улучшить работу подхода в отдельных случаях.
\end{abstract}

\begin{IEEEkeywords}
полуструктурированные данные, JSON, колоночные СУБД,
вертикальная декомпозиция, MVCC, ClickHouse
\end{IEEEkeywords}


\section{Введение}
\label{sec:intro}

Полуструктурированные данные в формате JSON стали стандартом для
обмена информацией в современных системах: журналы событий,
телеметрия, API-ответы, документы IoT. Объёмы таких данных растут, а
аналитические запросы по ним предъявляют требования, типичные для
OLAP: чтение значений отдельных полей миллиардов документов с минимумом
накладных расходов на десериализацию.

\textbf{Адресат и постановка задачи.} Настоящая работа адресована
\emph{разработчику СУБД}: автору реляционного движка, и призвана предложить
такой способ организации хранения JSON, который, во-первых,
обеспечивает аналитическое качество subcolumn-подхода ClickHouse, а
во-вторых, требует от разработчика минимальных изменений в уже
существующем физическом слое.

Канонические подходы, известные на сегодня, можно разделить на три
семейства.

\textit{Хранение JSON как неделимого значения.}
Документ сериализуется в бинарный формат (BSON в MongoDB, JSONB в
PostgreSQL, OSON в Oracle) и хранится в одной ячейке. Доступ к
отдельному полю требует разбора всего документа --- даже если из него
нужен один скаляр. Семантика максимально гибкая, эффективность
аналитических чтений минимальна.

\medskip
\textit{Entity-Attribute-Value (EAV) и его варианты.}
Документ раскладывается в строки общей таблицы
\texttt{(doc\_id, key, value)}. Поля разных типов либо приводятся к
строке, либо разносятся в несколько по-типовых таблиц (как в Argo3
\cite{argo}). Подход прост и не требует знания схемы, но плохо
сочетается с колоночным сжатием: значения одного и того же поля
оказываются физически разбросаны, а реконструкция документа требует
рекуррентного вызова операции JOIN.

\medskip
\textit{Колонка на каждый JSON-путь (subcolumn-подход).}
Современные колоночные движки (ClickHouse \cite{clickhouseJson},
BigQuery Capacitor \cite{bigqueryCapacitor}, SingleStore JSON over
Columnstore \cite{singlestoreJson}, Apache Doris VARIANT
\cite{dorisVariant}, Snowflake VARIANT \cite{snowflakeSigmod})
автоматически выводят схему по выборке документов и раскладывают
каждый встретившийся путь в собственный физический столбец. Доступ
к полю превращается в обычное чтение колонки. Данный подход
обеспечивает на порядки лучшую производительность аналитических
запросов по сравнению с двумя предыдущими, и в настоящей работе мы
ориентируемся именно на него.

\subsection{Проблема разработчика: subcolumn-подход требует
эффективного хранения колонок с большим числом NULL-значений}
\label{sec:problem}

ClickHouse мы берём за образец качества, которого хотим достичь.
Если же посмотреть на subcolumn-подход \emph{глазами разработчика},
желающего повторить его в своей СУБД, выясняется, что простая в
идее декомпозиция (каждый путь становится колонкой) на практике
сводится к одной конкретной инженерной задаче: \emph{уметь дёшево
хранить колонки, в которых подавляющее большинство значений
отсутствует}.

Документы реального JSON разреженные в схеме: сотни различимых
путей, но большинство встречается лишь в малой доле документов
(опциональные поля, поля, специфичные для разных типов событий или
версий данных, и так далее). Если каждый встречающийся путь
становится своей колонкой, то для каждого документа, в котором
этот путь не встретился, формально нужно хранить NULL. На корпусах
в миллиарды документов и тысячах путей это триллионы
NULL-значений, которые приходится чем-то выражать на физическом
уровне. Обычная плотная колонка с честным хранением NULL-маски в
этих условиях неприемлема ни по памяти, ни по скорости
сериализации/чтения. Без специализированного решения этого вопроса
весь смысл subcolumn-подхода теряется: его привлекательность
держится именно на том, что доступ к полю $p$ стоит как чтение
одной колонки, а не как разбор всего документа.

ClickHouse решает этот вопрос на физическом уровне:
страничный формат хранения, формат сериализации колонок,
compaction-логика и индексы изначально спроектированы с учётом
большого числа отсутствующих значений в subcolumn-ах. Эта
инфраструктура~--- значительный объём кода и часть архитектуры
движка, не сводимая к одной таблице или одному оператору.

Для разработчика, встраивающего поддержку JSON в \emph{существующий}
реляционный движок, эта инфраструктура с большой вероятностью
отсутствует. Воспроизведение её в чужом движке означает написание
собственной подсистемы хранения колонок с большим числом NULL-ов
и интеграцию её со всеми остальными слоями движка. Объём работы
сопоставим с написанием отдельной подсистемы хранения и плохо
совместим с целью ``минимально расширить существующий движок''.

Возникает практический вопрос: \emph{можно ли получить
аналитическое качество subcolumn-подхода, не реализуя эффективное
хранение колонок с большим числом NULL-ов на физическом уровне~---
а опираясь только на абстракции, уже имеющиеся в любой реляционной
СУБД (таблицы, индексы, MVCC)?} Цель настоящей работы --- показать,
что да, и что цена этого --- одна дополнительная физическая
операция при чтении (hash-gather по \texttt{row\_id}), которая на
современных аппаратных конфигурациях поглощается стандартными
оптимизациями.

\subsection{Вклад работы}

С позиции разработчика СУБД настоящая работа предоставляет
следующее.

\begin{enumerate}
  \item \emph{Схема хранения JSON-документов поверх стандартных
        абстракций реляционной СУБД} и алгоритмы выполнения над ней
        четырёх основных операций~--- чтения, вставки, удаления и
        изменения. От движка требуются только обычные таблицы,
        хеш-индексы, MVCC и стандартные физические операторы scan и
        hash-join; никаких изменений формата страниц, сериализации
        или compaction-логики не требуется. Это позволяет добавить
        эффективную поддержку JSON в широкий класс существующих
        движков без переделки физического слоя.

  \item \emph{Доказательство корректности схемы при параллельных
        запросах}. Сформулирован инвариант ``main~--- источник
        истины''; показано, что при заданном порядке записей внутри
        одной MVCC-транзакции (\textsc{insert}: side\,$\to$\,main,
        \textsc{delete}: main\,$\to$\,side) этот инвариант
        сохраняется на коммите каждой операции, что обеспечивает
        snapshot-isolation-консистентное поведение \textsc{select},
        \textsc{insert}, \textsc{delete} и \textsc{update} без
        введения каких-либо дополнительных механизмов синхронизации
        сверх уже имеющегося в движке MVCC-протокола.

  \item \emph{Три эвристики оптимизации базовой схемы}: отложенное
        удаление в side-таблицах с переносом работы в фоновую
        compaction; дублирующее хранение полного документа в main
        для эффективной точечной выборки по первичному ключу;
        битовая маска присутствия полей в main для сужения
        множества side-таблиц, к которым необходимо обращаться при
        удалении или поиске по \texttt{IS NOT NULL}. Каждая
        эвристика независима и может быть включена разработчиком
        выборочно в зависимости от профиля нагрузки.
\end{enumerate}

\section{Подход}
\label{sec:approach}

Ниже описывается схема в виде, удобном для разработчика-имплементатора:
указано, какие именно существующие примитивы движка задействуются на
каждом шаге, и какие новые сущности необходимо добавить.

\subsection{Декомпозиция документа}

Документ JSON раскладывается на корневые пути; каждый встречающийся
скалярный путь $p$ типа $\tau$ хранится в собственной
\emph{side-таблице}, которую будем обозначать
\texttt{side}($p$,\,$\tau$). Таблица \texttt{side}($p$,\,$\tau$)
имеет схему $(\texttt{row\_id}, \texttt{value})$, где
\texttt{row\_id}~--- ссылка на строку основной таблицы, которую
далее будем называть \texttt{main}, а \texttt{value}~--- значение
поля. На колонку \texttt{row\_id} в \texttt{side}($p$,\,$\tau$)
навешивается \textbf{хеш-индекс} для константного времени lookup
при реконструкции (см. подраздел~\ref{sec:scan}); без него
стоимость hash-gather была бы пропорциональна объёму side-таблицы.
В дальнейшем тексте мы используем эти обозначения~--- \texttt{main}
и \texttt{side}($p$,\,$\tau$)~--- единообразно. Реальное физическое
имя side-таблицы в имплементации может быть устроено любым
способом, лишь бы у каталога была возможность по паре $(p,\tau)$
получить ссылку на соответствующее отношение.

Сама таблица \texttt{main} в базовой схеме содержит только колонку
\texttt{row\_id} (первичный ключ); никакие данные документа в ней
не хранятся. Опциональные эвристики раздела~\ref{sec:heuristics}
могут добавить в \texttt{main} дополнительные служебные колонки,
но это не меняет основного устройства схемы.

Для документов, в которых одно и то же поле принимает значения
разных типов, заводится несколько side-таблиц~--- по одной на тип.
Это сознательно отказывает в обобщённой колонке Variant: вместо
дискриминатора на каждой строке, как в ClickHouse Dynamic, мы
получаем строгие узкие таблицы с гомогенным значением, для которых
сжатие и фильтрация работают без накладных расходов на ветвление по
типу. С точки зрения разработчика это означает, что не нужно
реализовывать тип-дискриминатор и его поддержку во всех физических
операторах: каждая side-таблица имеет обычный фиксированный тип
колонки \texttt{value}, такой же, как любая другая таблица в
движке.

\begin{figure}[h]
\centering
\begin{lstlisting}[caption={Логический разрез данных. Документы
$\{$user.id=42, user.name="ann", flag=true$\}$ и
$\{$user.id=7, level=3$\}$.},label={fig:layout}]
       main(row_id)
       +--------+
       | row_id |
       +--------+
       |   1    |
       |   2    |
       +--------+

       side(user.id, int)(row_id, value)
       +--------+-------+        [HASH INDEX on row_id]
       | row_id | value |
       +--------+-------+
       |   1    |   42  |
       |   2    |    7  |
       +--------+-------+

       side(user.name, str)(row_id, value)
       +--------+-------+        [HASH INDEX on row_id]
       | row_id | value |
       +--------+-------+
       |   1    | "ann" |
       +--------+-------+

       side(flag, bool)(row_id, value)
       +--------+-------+        [HASH INDEX on row_id]
       | row_id | value |
       +--------+-------+
       |   1    | true  |
       +--------+-------+

       side(level, int)(row_id, value)
       +--------+-------+        [HASH INDEX on row_id]
       | row_id | value |
       +--------+-------+
       |   2    |   3   |
       +--------+-------+
\end{lstlisting}
\end{figure}

В отличие от subcolumn-подхода, NULL-ов нет физически: отсутствие
поля у документа выражается отсутствием строки в соответствующей
side-таблице. Следствие~--- объём данных линейно зависит от числа
\emph{реально} присутствующих значений, а не от
$|\text{документы}| \times |\text{путей в схеме}|$.

\subsection{Гибридное сканирование}
\label{sec:scan}

Reconstruction документа выполняется специализированным физическим
оператором \texttt{scan\_computing\_table}. Логически он
эквивалентен выражению

\[
\texttt{main} \;\text{LEFT JOIN}\; \texttt{side}(p_1,\tau_1)
\;\text{LEFT JOIN}\;\dots\;\text{LEFT JOIN}\;
\texttt{side}(p_k,\tau_k)
\]
по \texttt{row\_id}, где $(p_1,\tau_1),\dots,(p_k,\tau_k)$~---
проецируемые пары путь/тип после column pruning. Физически
выполнение устроено так:

\begin{enumerate}
  \item Параллельно стартуют сканы \texttt{main} и всех нужных
        side-таблиц \texttt{side}($p_i$,\,$\tau_i$).
  \item Чанк \texttt{main} собирается полностью (включая
        \texttt{row\_id}).
  \item Для каждой пары $(p_i,\tau_i)$ запускается hash-gather: по
        множеству \texttt{row\_id}-ов чанка выполняется точечный
        lookup в хеш-индексе \texttt{side}($p_i$,\,$\tau_i$).
        Найденные значения раскладываются в вектор-колонку,
        остальные позиции остаются NULL.
\end{enumerate}

\begin{figure}[h]
\centering
\begin{lstlisting}[caption={Схема исполнения
\texttt{scan\_computing\_table}.},label={fig:scan}]
              +-------------------+
              |   scan main       |---+
              |   chunk: row_ids  |   |
              +-------------------+   |
                                      v
   +------------------+   +-----------------+   +------------------+
   | scan             |   | hash-gather by  |   | scan             |
   | side(p_1, t_1)   |-->| row_id from     |<--| side(p_k, t_k)   |
   +------------------+   | each side       |   +------------------+
   [HASH IDX]             +-----------------+   [HASH IDX]
                                   |
                                   v
                    chunk(main + p_1 + ... + p_k)
\end{lstlisting}
\end{figure}

Column pruning делает оператор естественно оптимизируемым: если
запрос проецирует только два поля, гибридный скан касается только
двух side-таблиц независимо от размера схемы документа.

\section{Корректность параллельных операций}
\label{sec:concurrency}

Хранилище состоит из множества физически разных таблиц. Любое
изменение документа атомарно затрагивает несколько из них, что
ставит классический вопрос согласованности при параллельных
читателях и писателях. Мы опираемся на действующий MVCC-протокол
движка и единственный дополнительный инвариант.

\paragraph{Инвариант MAIN.}
\emph{Если строка с идентификатором $r$ видима в \texttt{main} с
\texttt{commit\_id} $c$, то для любой пары $(p,\tau)$, по которой у
документа имеется значение типа $\tau$ по пути $p$, в
\texttt{side}($p$,\,$\tau$) существует видимая с
\texttt{commit\_id} $c'\le c$ строка $(r, v)$, и $v$~---
действительное значение поля.}

Иначе говоря, к моменту, когда строка появляется в main с
коммитом $c$, все её side-записи уже видимы с не более поздним
коммитом. Поддерживать это~--- задача процедур \textsc{insert},
\textsc{update} и \textsc{delete}.

\subsection{\textsc{Insert}}

Алгоритм multi-collection \textsc{insert}-а под одной транзакцией:

\begin{enumerate}
  \item Выделить $r$ как новый \texttt{row\_id} (атомный счётчик
        под локом таблицы или DB-wide-счётчик).
  \item Для каждого присутствующего в документе поля $(p,\tau,v)$
        вставить $(r, v)$ в соответствующую \texttt{side}($p$,\,$\tau$)
        со своим \texttt{txn\_id}; если такой таблицы ещё нет в
        каталоге, она создаётся.
  \item Только после всех side-вставок вставить строку в main
        со своим \texttt{txn\_id}.
  \item Зафиксировать транзакцию: единый \texttt{commit\_id}
        для всех физических вставок.
\end{enumerate}

Поскольку до момента commit ни одна из записей не видна другим
транзакциям (по правилам snapshot isolation), читатель либо вообще
не видит ни main, ни side, либо видит и main, и все side
одновременно. Промежуточное состояние, при котором main существует
без какой-либо из side, наружу не выходит. Инвариант MAIN
удовлетворяется.

\subsection{\textsc{Delete}}

Зеркальный порядок:

\begin{enumerate}
  \item Пометить строку main как удалённую (записать
        \texttt{delete\_txn\_id}).
  \item Пометить соответствующие строки во всех side-таблицах.
  \item Зафиксировать транзакцию.
\end{enumerate}

Обратный порядок обязателен: пока main видна как живая, читатель
имеет право обратиться в side. Если бы сначала удалялись side,
существовало бы временное окно, в котором main говорит ``поле
было'', а side говорит ``нет''~--- инвариант нарушился бы.

\subsection{\textsc{Update}}

\textsc{Update} раскладывается на пару \textsc{delete} +
\textsc{insert} с новым \texttt{row\_id}, либо на in-place
\textsc{update} side-таблиц с сохранением \texttt{row\_id}. В
обоих случаях операция выполняется под одним \texttt{txn\_id}; при
in-place варианте side-обновления применяются до main-обновления,
по тому же принципу, что и \textsc{insert}.

\subsection{Параллельность без сериализации запросов}

Существенно, что описанная корректность \emph{не} опирается на
факт ``одна коллекция $\to$ один исполнитель''. В реализациях
поверх акторной модели запросы на одну таблицу обычно сериализуются
маршрутизацией по хешу имени, что эквивалентно глобальному
сериализующему семафору. Эту маршрутизацию хочется уметь
ослаблять. Описанная схема к этому готова: безопасность
параллельных DML обеспечивается порядком записи (side$\to$main или
main$\to$side) под одной транзакцией и MVCC-протоколом, а не
эксклюзивным владением таблицей.

\section{Эвристики оптимизации}
\label{sec:heuristics}

Ниже описаны три необязательные эвристики, каждая из которых может
быть включена разработчиком независимо в зависимости от того, какие
шаблоны нагрузки наиболее важны для целевой СУБД. Эвристики не
меняют корректности базовой схемы и не вводят новых сущностей в
движке~--- только дополнительные колонки в main-таблице и
дополнительные шаги в существующих физических операторах.

\subsection{Отложенное удаление при compaction}

В описанной выше процедуре \textsc{delete} выполняются две группы
действий: проставление tombstone в main и проставление tombstone в
side-таблицах. С точки зрения корректности на пути чтения~---
важна только tombstone в main: если main говорит, что строка
$r$ удалена для снимка $s$, наличие живых записей в
\texttt{side}($p$,\,$\tau$) с \texttt{row\_id}\,$=r$ не влияет на
результат, потому что
гибридный скан фильтрует side по \texttt{row\_id}-ам, прошедшим
фильтр main, а не наоборот.

Это позволяет сделать удаление в side-таблицах \emph{отложенным}.
Процедура удаления документа сводится к одному действию~---
tombstone в main. Реальное удаление строк из side-таблиц
производится фоновой compaction-задачей, обходящей side-таблицы
и удаляющей строки, для которых \texttt{row\_id} помечен как
удалённый в main или вовсе отсутствует.

Преимущества:
\begin{itemize}
  \item \textsc{delete} становится O(1) по числу
        side-таблиц~--- работа с одной таблицей независимо от
        ширины схемы документа;
  \item compaction агрегирует множество удалений в один проход и
        работает в стиле LSM, без оверхеда на per-row I/O;
  \item конкурентный \textsc{delete} становится тривиально
        параллельным: транзакции на разные документы вообще не
        пересекаются по таблицам.
\end{itemize}

Цена~--- временное хранение мёртвых строк в side-таблицах до
compaction. Пока compaction не отработал, объём side-таблицы
завышен; колонка \texttt{value} хранит значения, которые никогда
не будут прочитаны. На рабочих нагрузках с подавляющей долей чтения
этот overhead ограничен соотношением частоты delete к частоте
compaction.

\subsection{Дублирующее хранение документа в main}

При типичной OLAP-нагрузке гибридный скан превосходит чтение
сериализованного документа благодаря column pruning. Однако
существует противоположный сценарий: точечный \textsc{select} по
первичному ключу, который возвращает один документ целиком (`OLTP
get-by-id'). В нём гибридный скан вынужден сделать $k$ hash-lookup
по числу путей в документе, что хуже единственного чтения
сериализованного JSON-блоба.

Эвристика: в main вводится дополнительная скрытая колонка
\texttt{\_raw} типа BYTEA, содержащая сериализованный документ в
компактном бинарном формате (BSON-подобный, MessagePack, или
internal binary representation). При построении плана для точечного
\textsc{select} оптимизатор выбирает чтение этой колонки и
декодирование вместо гибридного скана. Для аналитических запросов с
проекцией одного-двух полей выбирается классический гибридный скан.

Стоимость~--- ~$2\times$ память (документ хранится и в side, и в
main). Преимущество~--- сохраняется оптимальная стоимость как для
аналитических запросов, так и для точечных выборок. Дополнительный
плюс: \texttt{\_raw} позволяет восстановить полный документ без
обращения к side-таблицам, что упрощает резервное копирование и
репликацию.


\subsection{Битовая маска присутствия полей}

Третья эвристика обращается к проблеме \textsc{delete}-производи\-тель\-ности
в её строгом~--- не-отложенном~--- варианте.

Без дополнительной информации операция \textsc{delete} должна
проверить все $K$ side-таблиц схемы: даже если документ
содержит только два поля, мы не знаем заранее, в каких именно из
$K$ возможных side-таблиц лежат его записи. Соответственно,
\textsc{delete} выполняет $K$ обращений к side-таблицам, из
которых $K-k$ оказываются холостыми.

Эвристика: в main добавляется колонка \texttt{\_mask} типа
BITVECTOR фиксированной длины $K$ (или динамически-расширяемой
длины с per-таблицей мэппингом ``номер бита $\to$ имя пути''). Бит
$i$ установлен тогда и только тогда, когда в документе присутствует
значение по $i$-му пути. При \textsc{delete} запрос читает
\texttt{\_mask}, и затем обращается только к тем side-таблицам, для
которых соответствующий бит равен 1.

В сочетании с эвристикой отложенного удаления (раздел 5.1) маска
становится полезной не на пути \textsc{delete}, а на пути
compaction: фоновая задача может сужать множество side-таблиц,
которые надо проверить на наличие удалённых строк.

Кроме того, маска полезна для оптимизатора: на этапе планирования
запроса \texttt{SELECT path FROM t WHERE path IS NOT NULL} можно
заменить join со \texttt{side}($p$,\,$\tau$) на быстрое чтение
бита маски в \texttt{main},
если для конкретной строки явный value не нужен.

\section{Связанные работы}
\label{sec:related}

Идея вертикальной декомпозиции восходит к Decomposition Storage
Model \cite{dsm1985}, в которой каждый атрибут хранится как
бинарное отношение $(\textit{surrogate}, \textit{value})$. MonetDB
\cite{monetdb} реализует эту идею в чистом виде: каждая колонка
есть BAT (Binary Association Table). Наша работа применяет эту же
декомпозицию к полям JSON, но с принципиальной разницей в семантике
хранения: в side-таблицу попадают только реально присутствующие в
документах значения, тогда как BAT в MonetDB описывает плотный
набор пар $(\textit{oid}, \textit{value})$ для всех документов.

Argo \cite{argo} (Chasseur, Li, Patel, WebDB 2013) предлагает
хранить JSON в трёх таблицах по типам значений
($(\textit{docID}, \textit{key}, \textit{value})$). Это
ближайший по структуре подход. Отличия: Argo разлагает по типам, а
не по путям; ключ-путь хранится как значение колонки, а не как имя
таблицы (что лишает оптимизатор статической информации о схеме);
main-таблица отсутствует; колоночного слоя нет; вопрос корректности
параллельного DML не рассматривается. Argo проектировался как
OLTP-обёртка над RDBMS, не как аналитическое решение.

Промышленные subcolumn-движки --- ClickHouse JSON
\cite{clickhouseJson,clickhouseShared}, BigQuery Capacitor
\cite{bigqueryCapacitor}, SingleStore JSON over Columnstore
\cite{singlestoreJson}, Apache Doris VARIANT \cite{dorisVariant},
Snowflake VARIANT \cite{snowflakeSigmod}, Vertica Flex Tables
\cite{verticaFlex} --- получают аналогичный эффект, но требуют
специализированной поддержки разреженных колонок на физическом
уровне (см. раздел~\ref{sec:problem}). Наша работа отличается
именно тем, что обходится только стандартными реляционными
абстракциями.

Dremel \cite{dremel} вводит модель repetition/definition levels для
вложенных структур; SingleStore наследует эту модель. Мы
сознательно не используем repetition/definition levels:
вложенность раскладывается на отдельные пути, что делает массивы
менее естественно представимыми, но избавляет от необходимости
поддерживать специальный формат на нижнем уровне.

\section{Заключение}

Мы предложили \emph{рецепт для разработчика реляционной СУБД},
позволяющий добавить в движок эффективную поддержку JSON, не
переписывая физический слой хранения. Рецепт использует только
стандартные абстракции, доступные практически в любом реляционном
движке: обычные таблицы, хеш-индексы, MVCC, физические операторы
scan/join. Идейно подход повторяет subcolumn-разложение современных
колоночных движков (ClickHouse, BigQuery, SingleStore, Doris,
Snowflake), но переносит сложность с физического слоя в реляционный:
вместо формата с поддержкой большого числа NULL-ов мы используем
отсутствие строки в side-таблице, и вместо специализированного
формата для разделяемых редких путей --- множество узких таблиц.
Корректность параллельного DML обеспечена инвариантом
``main~--- источник истины'' и упорядочением записей внутри одной
MVCC-транзакции~--- без введения новых механизмов синхронизации сверх
уже имеющегося в движке протокола.

Предложенные три эвристики покрывают известные слабости базовой
схемы: \textsc{delete}-производительность,
точечную выборку и поиск присутствия поля.

Подход реализован в колоночной СУБД Otterbrix; экспериментальная
валидация на бенчмарке JSONBench подтверждает конкурентоспособность
относительно специализированных промышленных решений. Основным
направлением дальнейших исследований является адаптивная битовая
маска присутствия с экономичным кодированием для широких схем
документов и автоматическое продвижение часто используемых путей в
обычные колонки \texttt{main} по статистике обращений.

\bibliographystyle{IEEEtran}
\bibliography{references}

\end{document}
